\clearpage
\fancyhead[L]{\small EFFEMERIDI ASTRONOMICHE DI MILANO PER L'ANNO 1817.}
\hypertarget{native-606-lower}{}
\begin{center}\small Göttingische gelehrte Anzeigen. No. 144. Pages 1433..1438. 8 September 1817.\\[3pt]\rule{.72\linewidth}{.3pt}\end{center}

{\setlength{\rightskip}{0pt plus 3em}\emph{Effemeridi astronomiche di Milano per l'anno \textup{1817} calcolate da} \textsc{Francesco Carlini} \emph{ed} \textsc{Enrico Brambilla}. \emph{Con Appendice.} From the Imperial-Royal Printing House. Milan, 1816. 108 and 116 pages in octavo.\par}

\begingroup\brokenpenalty=10000\relax
Since the arrangement of the calendar in this volume and the accompanying supplements are exactly the same as in the preceding volume, this time we have only the papers in the appendix to report on. In the first, Mr.\ \textsc{Oriani} gives us the list of declinations of 40 stars drawn from the observations made with \textsc{Reichenbach}’s three-foot repeating circle and printed in earlier volumes of the ephemerides. The accompanying comparison with the values reported by \textsc{Maskelyne}, \textsc{Piazzi}, and \textsc{Pond} is noteworthy because of the instruments used by those astronomers. As regards the attainable accuracy, Mr.\ \textsc{Oriani} considers \textsc{Reichenbach}’s three-foot repeating circle roughly equal to \textsc{Troughton}’s five-foot mural circle at Greenwich. As for ease of observing, he likewise grants the latter instrument no advantage, on the grounds that the extra time required for a fourfold repetition is made up for by the greater ease of reading the repeating circle. Of course, this leaves out of account the time required to reduce the observations made off the meridian. --- The second paper, by \textsc{Ottaviano Fabrizio Mossotti}, a new analysis of the problem of determining the orbits of celestial bodies, deserves a somewhat detailed account because of the importance of the subject. The entire treatment is purely analytical and is carried out with skill. Yet the great number of symbols makes it somewhat harder to follow the thread, see the treatment as a whole, and grasp clearly what really matters. The essential point is as follows. If half the parameter of the orbit to be determined and half that of the Earth’s orbit are denoted by $p$ and $P$, the products of $\BareRoot{p}$ \SourceLeaf{607}with the cosines of the inclination of the plane of the unknown orbit to three fixed arbitrary planes by $c$, $c'$, $c''$, and likewise the products of $\BareRoot{P}$ with the cosines of the inclinations of the Earth’s orbit to those same fundamental planes by $C$, $C'$, $C''$, then Mr.\ \textsc{Mossotti} finds three linear equations of the form
\[
A(c-C)+A'(c'-C')+A''(c''-C'') = 0
\]
where the coefficients are known quantities. These equations, however, are only approximately correct insofar as the time intervals are not infinitely small, and are really based solely on the fact that the three segments formed when the three positions are joined pairwise by chords are very nearly proportional to the cubes of the time intervals, both in the unknown orbit and in the Earth’s orbit. Precisely because those equations are only approximately true, it is impermissible to combine them in such a way that one of the quantities $c-C$, $c'-C'$, $c''-C''$ would be eliminated. Rather, it can be shown that all three equations may properly count as only one, and Mr.\ \textsc{Massotti} accordingly retains only one of them. Since this equation alone can take us no further, Mr.\ \textsc{Massotti} enlists a fourth observation, which, when combined with two of the previous observations, yields a quite similar equation. Combining it with the equation deduced from the first three observations is indeed legitimate, and this combination is used to determine the ratio of the three quantities $c-C$, $c'-C'$, $c''-C''$ to one another, so that the results appear in this form
\[
c'-C' = M(c-C), \quad c''-C'' = N(c-C)
\]
By using the principle stated above, Mr.\ \textsc{Massotti} then also finds the ratio of the three distances of the celestial body from the Earth to one of the three quantities $c-C$, $c'-C'$, $c''-C''$, and finally the absolute values of these three quantities themselves. Since $C$, $C'$, $C''$ are known independently, the values of $c$, $c'$, $c''$ follow, and with them $p = cc+c'c'+c''c''$. From the expressions
\[
\frac{c}{\BareRoot{p}}, \quad \frac{c'}{\BareRoot{p}}, \quad \frac{c''}{\BareRoot{p}}
\]
we also obtain the quantities on which the orientation of the celestial body’s orbit depends. The further determination of the remaining elements rests on known theorems and presents nothing unusual.

If we are now to give our judgment of this new method here, we first acknowledge with pleasure that the aforementioned linear equation between $c-C$, $c'-C'$, $c''-C''$ is noteworthy and that it is interesting to see the possibility of an entirely linear determination of an unknown orbit worked out. A different question, however, is whether such a method is to be recommended for actual practical use. It is clear that not everything contained in the theory of motion in conic sections has been used here and that not everything was meant to be used. In particular, to put it in the spirit of a method of approximation, the condition is entirely ignored that the quotients obtained by dividing those segments by the cubes of the time intervals must be very nearly equal to the quantity $\dfrac{k^3}{12r^3}$, where $k$ is to be taken in the sense of the Theoria Motus Corporum Coelestium and $r$ denotes the distance of the celestial body from the Sun. Hence the necessity of resorting to a fourth observation in a problem where three observations already contain the complete solution. Even apart from the fact that this is theoretically inappropriate, the procedure becomes very precarious in practice because here the determination of the orbit actually rests on quantities of the \emph{third} order. The inevitable observational errors must therefore exert a disproportionate influence on the results, an influence that will be all the greater because the method, as a mere method of approximation, can allow only short \SourceLeaf{608}time intervals to be used. This method will therefore generally be able to give only results of little reliability, where making full use of three observations can already lead to a very close approximation.

Finally, we must also note that, if one chooses to start solely from that principle of the proportionality of the segments to the cubes of the time intervals, one can achieve in a simpler way the same result of a purely linear determination of the orbit from four observations. For that principle readily gives, from three observations, an equation, likewise linear, between two distances from the Earth, and in the same way one finds a second, similar equation between the same distances by combining the two corresponding observations with a fourth observation. Elimination then gives those distances themselves, just as accurately as the values they take when \textsc{Mossotti}’s method is applied, or rather, in fact, the very same values. As soon as these distances are known, determining the entire orbit presents no difficulty, as is well known. We must reserve the further development of this subject for another occasion, as this is naturally not the place for it. We must note only one more thing. The 80-page paper printed in the present annual volume of the ephemerides contains only the theoretical part of \textsc{Mossotti}’s method, and so a promised application to \textsc{Halley}’s comet will probably follow in the next volume. --- The present volume closes with a short article by \textsc{Carlini}, in which some errors in \textsc{Delambre}’s tables for the satellites of Jupiter are corrected.
\par\clearpage\endgroup
